\section{Tempo e spazio}

\subsection{Complessitá di tempo}
$M(x)$ richiede tempo $t$ sse:
\[
\configurazionestart{x} \xrightarrow{M^t} \configurazione{q_H}{w}{u}
\]
Dove $q_H \in \{h, "yes", "no"\}. $
\begin{itemize}
	\item {Una MdT deterministica $M$ {opera} in tempo $f(n)$ sse $\forall x \in \Sigma^*$ $M(x)$ richiede tempo $t \leq f(n)$} 
	\item {Una MdT non deterministica $N$ {opera} in tempo $f(n)$ sse $\forall x \in \Sigma^*$ tutti i run della computazione di $N(x)$ terminano in tempo $t \leq f(n)$} 
\end{itemize}


\subsection{Complessitá di spazio}
Si potrebbe considerare la somma della lunghezza dei $k$ nastri: $\sum_{i=1}^{k} |w_i u_i|$ \\
Ma possiamo scegliere anche la lunghezza del nastro piú lungo: $\max_{i \in \{1, \dots, k\}} |w_i u_i|$. \\
Queste due quantitá sono asintoticamente equivalenti:\\
$\max_{i \in \{1, \dots, k\}} |w_i u_i| \leq  \sum_{i=1}^{k} |w_i u_i| \leq k \cdot \max_{i \in \{1, \dots, k\}} |w_i u_i|$  \\

\textbf{Attenzione:} con queste definizioni naive, lo spazio occupato dall'input verrá sempre contato.
Questo vuol dire che lo spazio é sempre limitato inferiormente da $n$, anche quando non usiamo variabili ausiliarie. \\
A noi interessa considerare lo \textit{spazio ausiliario} utilizzato. \\
Per farlo, occorre considerare le macchine di Turing con input e output e contare solo i nastri interni.

\subsubsection{Complessitá di spazio ausiliario nelle MdT con input e output}
$M(x)$ richiede spazio $z$ sse: \\
Se:
\[
\configurazionestart{x} \xrightarrow{M^*} \configurazione{q_H}{w}{u}
\]
Dove $q_H \in \{h, "yes", "no"\} \quad$ e {$\quad z = \sum_{i=2}^{k-1} |w_i u_i|$}

\begin{itemize}
\item {Una MdT deterministica $M$ {opera} in spazio $g(n)$ sse $\forall x \in \Sigma^*$ $M(x)$ richiede spazio $z \leq g(n)$}
\item {Una MdT non deterministica $N$ {opera} in spazio $g(n)$ sse $\forall x \in \Sigma^*$ tutti i run della computazione di $N(x)$ terminano in spazio $z \leq g(n)$}
\end{itemize}


\subsection{Classi di complessitá}
\subsubsection{TIME}
Un linguaggio $L$ appartiene alla classe $\TIME(f(n))$ se esiste una macchina di Turing $M$ che decide $L$ e opera in tempo $f(n)$.
\subsubsection{NTIME}
Un linguaggio $L$ appartiene alla classe $\NTIME(f(n))$ se esiste una macchina di Turing non deterministica $N$ che decide $L$ e opera in tempo $f(n)$.
\subsubsection{SPACE}
Un linguaggio $L$ appartiene alla classe $\SPACE(f(n))$ se esiste una macchina di Turing $M$ che decide $L$ e opera in spazio $f(n)$.
\subsubsection{NSPACE}
Un linguaggio $L$ appartiene alla classe $\NSPACE(f(n))$ se esiste una macchina di Turing non deterministica $N$ che decide $L$ e opera in spazio $f(n)$.
\subsubsection{P}
$\P = \bigcup_{j>0} \TIME(n^j)$
\subsubsection{NP}
$\NP = \bigcup_{j>0} \NTIME(n^j)$
\subsubsection{EXP}
$\EXP = \bigcup_{j>0} \TIME(2^{n^j})$
\subsubsection{NEXP}
$\NEXP = \bigcup_{j>0} \NTIME(2^{n^j})$
\subsubsection{L}
$\L = \bigcup_{j>0} \SPACE(\log(n)^j)$
\subsubsection{PSPACE}
$\PSPACE = \bigcup_{j>0} \SPACE(n^j)$

\subsection{Classi di complessitá complementari}
Per ogni classe di complessitá $C$ definiamo la classe complementare $\co C$ come l'insieme dei linguaggi complementari a quelli di $C$. \\
$\co C = \{ \overline{L} : L \in C \}$


\subsection{Teorema: il tempo é limitato inferiormente dall'input}
\label{teorema:tempo_limitato_inferiormente_dall_input}
{$f(n) = \Omega(n)$} \\
O, equivalentemente: \\
{$n = O(f(n))$}
\begin{proof}
Se la codifica dell'input é \textit{ragionevole}, allora la macchina di Turing deve leggere tutto l'input prima di poter decidere se accettarlo o rifiutarlo.
\end{proof}


\subsection{Teorema: lo spazio é limitato superiormente dal tempo}
\label{teorema:spazio_limitato_superiormente_dal_tempo}
{$g(n) = O(f(n))$}
\begin{proof}
Supponiamo che una macchina $M$ a 1 nastro operi in tempo $f(n)$ e che a ogni passo scriva sul nastro e sposti la testina a destra. \\
Allora, dopo $f(n)$ passi, la testina avrá scritto esattamente $g(n) = f(n)$ simboli sul nastro. \\
Qualsiasi altra macchina che operi in tempo $f(n)$ avrá eseguito dei passi sovrascrivendo celle giá utilizzate. In questo caso l'utilizzo di spazio é inferiore a quello di $M$.
\end{proof}

\subsubsection{Corollario: $\TIME(f(n)) \subseteq \SPACE(f(n))$}
Un linguaggio che puó essere deciso in tempo $f(n)$ puó essere deciso anche in spazio $f(n)$. \\
Viceversa, non é possibile risolvere un problema in meno tempo di quanto serve a occupare lo spazio richiesto.
\footnote{Sappiamo che strutture dati come le hash table o tecniche di programmazione dinamica scaricano sullo spazio la complessitá di tempo. Ma dal punto di vista della teoria della complessitá, queste tecniche cambiano il problema alla radice.}

\subsubsection{Corollario: $\P \subseteq \PSPACE$}

\subsubsection{Raffinamento: $\TIME(f(n)) \subseteq \SPACE{ ( \frac{f(n)}{\log(f(n))} ) }$}
\label{raffinamento:time_space}
Vedere appendice (\ref{proof:raffinamento_time_space})


